% =====================================================================
\chapter{Números complejos en forma binómica}
\label{tema:5}
% =====================================================================

Los números complejos amplían los reales con una raíz cuadrada de $-1$. Con ellos toda ecuación polinómica tiene solución, y las operaciones de giro y homotecia del plano se escriben como productos. Este tema construye $\C$, demuestra las propiedades del conjugado y del módulo y resuelve ecuaciones de segundo grado con coeficientes complejos. Se sigue la §2.4 de~\cite{osat}, la sección~5 de las notas de 18.03 del MIT~\cite{mit1803} y la §1.1 de las notas de «Vectors and Matrices» de Cambridge~\cite{cambvm}.

\section{El cuerpo de los números complejos}

\begin{definicion}[Números complejos]\label{def:complejos}
Un \emph{número complejo} es una expresión $z=a+bi$ con $a,b\in\R$; se llaman \emph{parte real} $\Real z=a$ y \emph{parte imaginaria} $\Imag z=b$. Dos complejos son iguales si tienen iguales las partes reales y las imaginarias. Con $z=a+bi$ y $w=c+di$ se definen
\[
z+w=(a+c)+(b+d)i,\qquad zw=(ac-bd)+(ad+bc)i .
\]
El conjunto de los números complejos se denota $\C$. Se identifica el real $a$ con $a+0i$, y se escribe $i=0+1i$.
\end{definicion}

La regla del producto es la que resulta de operar como con polinomios en $i$ y sustituir $i^{2}=-1$: en efecto, con la definición, $i\cdot i=(0\cdot0-1\cdot1)+(0\cdot1+1\cdot0)i=-1$. Geométricamente, $a+bi$ se representa como el punto $(a,b)$ del \emph{plano complejo}, con el eje real horizontal y el imaginario vertical (\cref{fig:argand}).

\begin{teorema}\label{teo:cuerpo}
Con estas operaciones, $\C$ es un cuerpo: la suma y el producto son asociativos y conmutativos, el producto es distributivo respecto de la suma, $0$ y $1$ son los elementos neutros, y todo $z=a+bi\ne0$ tiene inverso
\[
z^{-1}=\frac{a}{a^{2}+b^{2}}-\frac{b}{a^{2}+b^{2}}\,i .
\]
\end{teorema}

\addtocounter{definicion}{1}%

\section{El conjugado}

\begin{definicion}[Conjugado]
El \emph{conjugado} de $z=a+bi$ es $\bar z=a-bi$. Geométricamente, es la reflexión de $z$ en el eje real.
\end{definicion}

\begin{proposicion}[Propiedades del conjugado]\label{prop:conjugado}
Para todos $z,w\in\C$:
\begin{enumerate}
  \item $\overline{z+w}=\bar z+\bar w$, \quad $\overline{zw}=\bar z\,\bar w$, \quad $\overline{z/w}=\bar z/\bar w$ si $w\ne0$, \quad $\bar{\bar z}=z$;
  \item $z+\bar z=2\Real z$ y $z-\bar z=2i\Imag z$;
  \item $z$ es real si y solo si $z=\bar z$;
  \item $z\bar z=(\Real z)^{2}+(\Imag z)^{2}$, que es un número real mayor o igual que~0.
\end{enumerate}
\end{proposicion}

La propiedad~(iv) da el método práctico para dividir: se multiplican numerador y denominador por el conjugado del denominador,
\[
\frac{z}{w}=\frac{z\,\bar w}{w\,\bar w}=\frac{z\,\bar w}{c^{2}+d^{2}} .
\]

\addtocounter{definicion}{1}%

\section{El módulo}

\begin{definicion}[Módulo]
El \emph{módulo} de $z=a+bi$ es $\abs{z}=\sqrt{a^{2}+b^{2}}=\sqrt{z\bar z}$, la distancia de $z$ al origen. La distancia entre $z$ y $w$ es $\abs{z-w}$.
\end{definicion}

\begin{proposicion}[Propiedades del módulo]\label{prop:modulo}
Para todos $z,w\in\C$:
\begin{enumerate}
  \item $\abs{z}\ge0$, con igualdad si y solo si $z=0$; \quad $\abs{\bar z}=\abs{z}$;
  \item $\abs{\Real z}\le\abs{z}$ y $\abs{\Imag z}\le\abs{z}$;
  \item $\abs{zw}=\abs{z}\abs{w}$ y, si $w\ne0$, $\abs{z/w}=\abs{z}/\abs{w}$;
  \item (desigualdad triangular) $\abs{z+w}\le\abs{z}+\abs{w}$;
  \item $\bigl|\abs{z}-\abs{w}\bigr|\le\abs{z-w}$.
\end{enumerate}
\end{proposicion}

\section{Ecuaciones polinómicas}

\begin{proposicion}[Raíces cuadradas]\label{prop:raiz-cuadrada}
Sea $w=a+bi$ con $b\ne0$. Las soluciones de $z^{2}=w$ son $z=\pm(x+yi)$ con
\[
x=\sqrt{\frac{\abs{w}+a}{2}},\qquad y=\operatorname{signo}(b)\sqrt{\frac{\abs{w}-a}{2}} .
\]
\end{proposicion}

\addtocounter{definicion}{1}%

\begin{proposicion}[Fórmula cuadrática en $\C$]\label{prop:cuadratica}
Sean $a,b,c\in\C$ con $a\ne0$, y sea $\delta\in\C$ con $\delta^{2}=b^{2}-4ac$. Las soluciones de $az^{2}+bz+c=0$ son $z=\dfrac{-b\pm\delta}{2a}$.
\end{proposicion}

\begin{proposicion}[Raíces conjugadas]\label{prop:conjugadas}
Si $p(z)=\sum_{k=0}^{n}c_kz^{k}$ tiene coeficientes reales, entonces $p(\bar z)=\overline{p(z)}$ para todo $z$. En particular, si $z$ es raíz de $p$, también lo es $\bar z$.
\end{proposicion}

\begin{teorema}[Teorema fundamental del álgebra]\label{teo:fundamental}
Todo polinomio de grado $n\ge1$ con coeficientes complejos tiene al menos una raíz en $\C$. En consecuencia, se factoriza como $c\prod_{k=1}^{n}(z-z_k)$ y tiene exactamente $n$ raíces contadas con multiplicidad.
\end{teorema}

La demostración requiere herramientas de análisis y se omite; véanse~\cite[§5.3]{mit1803} y~\cite{needham}. La consecuencia se obtiene por inducción sobre $n$, dividiendo por $z-z_1$ cada vez que se encuentra una raíz.

\section{El plano complejo y el producto escalar}

\addtocounter{figure}{1}%

Al identificar $\C$ con $\R^{2}$, la suma de complejos es la suma de vectores, y el producto por un real $\lambda$ es el producto de un vector por un escalar. El producto de dos complejos no tiene análogo en $\R^{2}$ y es el objeto del \cref{tema:6}. El conjugado permite recuperar el producto escalar y el área orientada:

\begin{proposicion}\label{prop:escalar-complejo}
Sean $z=a+bi$ y $w=c+di$, identificados con $u=(a,b)$ y $v=(c,d)$. Entonces
\[
\Real(z\bar w)=ac+bd=u\cdot v,\qquad \Imag(\bar z w)=ad-bc=\det\begin{pmatrix}a&c\\ b&d\end{pmatrix}.
\]
\end{proposicion}

\addtocounter{definicion}{1}%

\addtocounter{definicion}{1}%

